\documentclass[11pt,a4paper]{article}
\usepackage{amsmath,amssymb}
\usepackage{listings}
\usepackage{hyperref}
\usepackage[margin=1in]{geometry}
\usepackage{fancyhdr}

\pagestyle{fancy}
\fancyhf{}
\fancyhead[L]{Module 15: Game Development Projects}
\fancyhead[R]{C Programming Course}
\fancyfoot[C]{\thepage}

\lstset{
  language=C,
  basicstyle=\ttfamily\small,
  keywordstyle=\bfseries,
  commentstyle=\itshape,
  numbers=left,
  numberstyle=\tiny,
  frame=single,
  breaklines=true,
  tabsize=4
}

\title{Module 15: Game Development Projects}
\author{C Programming Course}
\date{}

\begin{document}
\maketitle
\thispagestyle{fancy}

\section{Console-Based Tic-Tac-Toe}
\begin{lstlisting}
#include <stdio.h>

static char board[3][3];

void init_board(void) {
    int cell = '1';
    for (int i = 0; i < 3; i++)
        for (int j = 0; j < 3; j++)
            board[i][j] = cell++;
}

void print_board(void) {
    for (int i = 0; i < 3; i++) {
        printf(" %c | %c | %c \n",
               board[i][0], board[i][1],
               board[i][2]);
        if (i < 2) printf("---+---+---\n");
    }
}
\end{lstlisting}

\section{Snake Game}
Key concepts for a console snake game:
\begin{itemize}
  \item Represent the snake as a linked list of coordinates.
  \item Use a 2D array for the game board.
  \item Poll for keyboard input to change direction.
  \item Detect collisions with walls and the snake body.
\end{itemize}

\section{Sudoku Solver}
A backtracking algorithm tries each digit 1--9 in empty cells and backtracks on
conflict.

\begin{lstlisting}
int is_safe(int grid[9][9], int row, int col,
            int num) {
    for (int i = 0; i < 9; i++) {
        if (grid[row][i] == num) return 0;
        if (grid[i][col] == num) return 0;
    }
    int sr = row - row % 3, sc = col - col % 3;
    for (int i = 0; i < 3; i++)
        for (int j = 0; j < 3; j++)
            if (grid[sr + i][sc + j] == num)
                return 0;
    return 1;
}
\end{lstlisting}

\section{Game Loop and State Management}
\begin{lstlisting}
#include <stdio.h>

typedef enum { MENU, PLAYING, GAME_OVER } State;

int main(void) {
    State state = MENU;
    int running = 1;
    while (running) {
        switch (state) {
            case MENU:
                printf("Press 1 to play\n");
                state = PLAYING;
                break;
            case PLAYING:
                printf("Game running...\n");
                state = GAME_OVER;
                break;
            case GAME_OVER:
                printf("Game Over!\n");
                running = 0;
                break;
        }
    }
    return 0;
}
\end{lstlisting}

A game loop typically consists of: process input, update state, render output.

\end{document}
